% Generated by build.mjs from ledger.mjs. Do not edit by hand.
\begin{longtable}{@{}p{1.25cm}p{9.3cm}P{2.5cm}l@{}}
\toprule ID & Requirement & Anchor & Class\\ \midrule \endhead
\bottomrule \endfoot
\leavevmode\phantomsection\label{YM-01}\textsf{YM-01} & \textbf{Target object.} For each admissible $G$ a quantum field theory $\mathcal T_G=(\mathcal H,U,\Omega,\mathcal D,\{\mathcal O_i\}_{i\in I})$ on Minkowski space $\mathbb R^{1,3}$ in the sense of Definition~\ref{def:wightman}. & \footnotesize S1 \S3--4 & \classbadge{NORMATIVE}\\[2pt]
\leavevmode\phantomsection\label{YM-02}\textsf{YM-02} & \textbf{All gauge groups.} The statement holds for \emph{every} compact simple gauge group. Constants (including $\Delta$) may depend on $G$; uniformity in $G$ is \emph{not} required. \emph{(Reading: \idref{I-2}.)} & \footnotesize S1 \S4 & \classbadge{NORMATIVE}\\[2pt]
\leavevmode\phantomsection\label{YM-03}\textsf{YM-03} & \textbf{Space-time $\mathbb R^4$.} The theory lives on $\mathbb R^4$: not on a lattice, not on a torus $T^4$, not in dimension $d<4$, and without any remaining UV or IR cut-off. & \footnotesize S1 \S4, \S6 & \classbadge{NORMATIVE}\\[2pt]
\leavevmode\phantomsection\label{YM-04}\textsf{YM-04} & \textbf{Non-triviality.} $\mathcal T_G$ is non-trivial. The term is not defined in S1; working definition in \idref{I-1}. \emph{(Reading: \idref{I-1}.)} & \footnotesize S1 \S4 & \classbadge{NORMATIVE}\\[2pt]
\leavevmode\phantomsection\label{YM-05}\textsf{YM-05} & \textbf{Hilbert space and symmetry.} A separable Hilbert space $\mathcal H$ with a strongly continuous unitary representation $U$ of the restricted Poincar\'e group $\mathcal P_+^\uparrow$ (or its covering group). & \footnotesize S1 \S3; S3 & \classbadge{NORMATIVE}\\[2pt]
\leavevmode\phantomsection\label{YM-06}\textsf{YM-06} & \textbf{Spectrum condition.} The joint spectrum of the generators $(H,\vec P)$ lies in the closed forward cone $\overline V_+$; in particular $H\ge 0$. & \footnotesize S1 \S3--4 & \classbadge{NORMATIVE}\\[2pt]
\leavevmode\phantomsection\label{YM-07}\textsf{YM-07} & \textbf{Vacuum.} A Poincar\'e-invariant vector $\Omega\in\mathcal H$, unique up to a phase. & \footnotesize S1 \S3 & \classbadge{NORMATIVE}\\[2pt]
\leavevmode\phantomsection\label{YM-08}\textsf{YM-08} & \textbf{Local fields.} Operator-valued tempered distributions $\mathcal O_i$ on a common dense $U$-invariant domain $\mathcal D\ni\Omega$, Poincar\'e-covariant, with $\Omega$ cyclic for the polynomial field algebra. & \footnotesize S1 \S3; S3 & \classbadge{NORMATIVE}\\[2pt]
\leavevmode\phantomsection\label{YM-09}\textsf{YM-09} & \textbf{Locality.} $[\mathcal O_i(f),\mathcal O_j(g)]=0$ on $\mathcal D$ whenever $\operatorname{supp}f$ and $\operatorname{supp}g$ are space-like separated. & \footnotesize S1 \S3 & \classbadge{NORMATIVE}\\[2pt]
\leavevmode\phantomsection\label{YM-10}\textsf{YM-10} & \textbf{Field content.} Local fields correspond to gauge-invariant local polynomials in $F$ and its covariant derivatives (e.g.\ $\operatorname{Tr}F_{ij}F_{kl}$). The correspondence is dimension-filtered, not a naive bijection. & \footnotesize S1 \S4 and fn.~1 & \classbadge{NORMATIVE}\\[2pt]
\leavevmode\phantomsection\label{YM-11}\textsf{YM-11} & \textbf{Short-distance identity.} Correlation functions agree at short distances with asymptotic freedom and perturbative renormalization theory. \emph{(Reading: \idref{I-4}.)} & \footnotesize S1 \S4 & \classbadge{NORMATIVE}\\[2pt]
\leavevmode\phantomsection\label{YM-12}\textsf{YM-12} & \textbf{Stress tensor and OPE.} Existence of a stress tensor and of an operator product expansion with the local singularities predicted by asymptotic freedom. & \footnotesize S1 \S4 & \classbadge{NORMATIVE}\\[2pt]
\leavevmode\phantomsection\label{YM-13}\textsf{YM-13} & \textbf{Mass gap.} $\exists\,\Delta>0:\ \sigma(H)\cap(0,\Delta)=\varnothing$. & \footnotesize S1 \S4 & \classbadge{NORMATIVE}\\[2pt]
\leavevmode\phantomsection\label{YM-14}\textsf{YM-14} & \textbf{Finite mass.} $m:=\sup\{\Delta\}<\infty$. Equivalently $H\neq 0$; with \idref{A-027} this excludes exactly $\mathcal H=\mathbb C\Omega$. & \footnotesize S1 \S4 & \classbadge{NORMATIVE}\\[2pt]
\leavevmode\phantomsection\label{YM-15}\textsf{YM-15} & \textbf{Axiomatic strength.} Axiomatic properties at least as strong as S3 and S4. On the Euclidean route: the OS axioms in the corrected form of \cite{OS75}, including the linear growth condition (\idref{A-001}, \idref{A-002}). \emph{(Reading: \idref{I-5}.)} & \footnotesize S1 \S4 & \classbadge{NORMATIVE}\\[2pt]
\leavevmode\phantomsection\label{YM-16}\textsf{YM-16} & \textbf{No weak existence.} Existence by compactness (subsequential limits) alone is excluded unless the axioms and the mass gap are established for the limit. & \footnotesize S1 fn.~2 & \classbadge{NORMATIVE}\\[2pt]
\leavevmode\phantomsection\label{YM-17}\textsf{YM-17} & \textbf{UV regulator removed.} Every UV regulator (lattice spacing $a$, momentum cut-off, \dots) is removed; the topology of convergence of the correlation functions is specified. & \footnotesize from YM-03 & \classbadge{DERIVED}\\[2pt]
\leavevmode\phantomsection\label{YM-18}\textsf{YM-18} & \textbf{IR regulator removed.} Every IR regulator (volume $L$, boundary conditions) is removed; dependence on boundary conditions is absent or controlled. & \footnotesize from YM-03 & \classbadge{DERIVED}\\[2pt]
\leavevmode\phantomsection\label{YM-19}\textsf{YM-19} & \textbf{Properties survive limits.} Reflection positivity, full Euclidean invariance $E(4)$ (restored from any lattice subgroup), locality, gauge invariance of observables and the gap bound hold for the \emph{limit}, not only along the approximation. & \footnotesize from YM-05--16 & \classbadge{DERIVED}\\[2pt]
\leavevmode\phantomsection\label{YM-20}\textsf{YM-20} & \textbf{Complete proof.} A complete mathematical solution: every step is proved, or reduced to a published theorem whose hypotheses are verified in the instance. & \footnotesize S2 \S5(a) & \classbadge{PROCEDURAL}\\[2pt]
\leavevmode\phantomsection\label{YM-21}\textsf{YM-21} & \textbf{Explicit correspondence.} The paper addresses the specific questions of the official description; a closely related question does not qualify. & \footnotesize S2 \S5(d) & \classbadge{PROCEDURAL}\\[2pt]
\leavevmode\phantomsection\label{YM-22}\textsf{YM-22} & \textbf{Self-containment.} CMI does not consider supplementary material submitted by the author: every essential step must be in the published work. & \footnotesize S2 \S6(d) & \classbadge{PROCEDURAL}\\[2pt]
\leavevmode\phantomsection\label{YM-23}\textsf{YM-23} & \textbf{Qualifying Outlet.} Publication in a refereed mathematics publication of worldwide repute with named editorial board, qualified editors, a published refereeing process and MathSciNet listing. CMI keeps no list and gives no advice. & \footnotesize S2 \S4(a), \S6 & \classbadge{PROCEDURAL}\\[2pt]
\leavevmode\phantomsection\label{YM-24}\textsf{YM-24} & \textbf{Time and acceptance.} At least two years since publication, survival of rigorous examination and general acceptance in the global mathematics community; only then may CMI decide on detailed consideration by a special advisory committee. & \footnotesize S2 \S4(b)--(c), \S7 & \classbadge{PROCEDURAL}\\[2pt]
\leavevmode\phantomsection\label{YM-25}\textsf{YM-25} & \textbf{Priority record.} CMI pays special attention to dependence on prior published insights: keep a contribution and priority ledger from day one. & \footnotesize S2 \S8(b) & \classbadge{GOVERNANCE}\\[2pt]
\end{longtable}
